% Môn: Vật lý
% Lớp: 11
% Chương: Chương I. Dao động
% Bài: L11C1 Bài 1. Dao động điều hòa
% Số câu: 143
% App đọc file này trực tiếp — sửa rồi Commit trên GitHub, bấm Đồng bộ GitHub .tex

% L11C1 Bài 1. Dao động điều hòa

% ===== Câu 11 =====
% ID: L11C1B1-04-DS
% Mức: TH
\dangbt{Mối liên hệ giữa dao động điều hòa và chuyển động tròn đều. Bài toán đếm số dao động}
\begin{ex}
% ID: L11C1B1-04-DS
% Mức: TH
Một vật dao động điều hòa có phương trình $x=4\cos(2\pi t+0)$ cm. Xác định tính đúng sai.
\choiceTF
{\True Biên độ bằng $4\,\text{cm}$.}
{\True Tần số góc bằng 2 $\pi$ rad/s.}
{\True Tần số bằng 2/ $2\,\text{Hz}$.}
{\True Chu kì bằng $1\,\text{s}$.}
\end{ex}

% ===== Câu 12 =====
% ID: L11C1B1-04-TLN
% Mức: TH
\begin{ex}
% ID: L11C1B1-04-TLN
% Mức: TH
Một vật thực hiện 12 dao động trong $6\,\text{s}$. Tần số dao động theo đơn vị Hz bằng bao nhiêu?
\shortans{2}
\end{ex}

% ===== Câu 18 =====
% ID: L11C1B1-06-TL
% Mức: TH
\begin{ex}
% ID: L11C1B1-06-TL
% Mức: TH
Một vật dao động theo phương trình $x=4\cos(2\pi t+\pi/3)$ cm. Xác định biên độ, tần số góc, chu kì, tần số và pha ban đầu.
\choiceTF

\end{ex}

% ===== Câu 22 =====
% ID: L11C1B1-08-TN
% Mức: TH
\begin{ex}
% ID: L11C1B1-08-TN
% Mức: TH
Một vật thực hiện dao động điều hòa trên một đoạn thẳng quỹ đạo dài $12$ cm. Biên độ dao động của vật là:
\choice
{$12$ cm}
{\True $6$ cm}
{$24$ cm}
{$3$ cm}
\loigiai{
• Quỹ đạo chuyển động của vật dao động điều hòa là một đoạn thẳng nối liền hai vị trí giới hạn xa nhất (biên âm và biên dương).
 
• Chiều dài quỹ đạo chuyển động $L$ liên hệ với biên độ dao động $A$ theo công thức: $L = 2A$.
 
• Từ số liệu đề bài cho $L = 12$ cm, ta tính được biên độ dao động là: $A = \dfrac{L}{2} = \dfrac{12}{2} = 6$ cm.
}
\end{ex}

% ===== Câu 25 =====
% ID: L11C1B1-09-TN
% Mức: TH
\begin{ex}
% ID: L11C1B1-09-TN
% Mức: TH
Một chất điểm chuyển động tròn đều trên đường tròn có bán kính $R = 10$ cm với tốc độ góc bằng $5$ rad/s. Hình chiếu của chất điểm này lên một đường kính nằm ngang dao động điều hòa với biên độ và tần số góc lần lượt là:
\choice
{$5$ cm; $10$ rad/s}
{\True $10$ cm; $5$ rad/s}
{$20$ cm; $5$ rad/s}
{$10$ cm; $10\pi$ rad/s}
\loigiai{
• Dựa trên mối liên hệ giữa chuyển động tròn đều và dao động điều hòa: Hình chiếu của một chất điểm chuyển động tròn đều lên một đường kính nằm trong mặt phẳng quỹ đạo dao động điều hòa.
 
• Biên độ dao động $A$ của hình chiếu bằng bán kính của quỹ đạo chuyển động tròn đều: $A = R = 10$ cm.
 
• Tần số góc $\omega$ của dao động điều hòa bằng tốc độ góc $\omega$ của chuyển động tròn đều: $\omega = 5$ rad/s.
}
\end{ex}

% ===== Câu 26 =====
% ID: L11C1B1-10-DS
% Mức: VD
\dangbt{Xác định chu kì, tần số và biên độ từ các dữ kiện quỹ đạo và thời gian}
\begin{ex}
% ID: L11C1B1-10-DS
% Mức: VD
Một chất điểm chuyển động tròn đều với tốc độ góc bằng $5\pi$ rad/s trên đường tròn có bán kính bằng $6$ cm. Hình chiếu của chất điểm này lên một đường kính nằm ngang dao động điều hòa. Khảo sát các thông số của hình chiếu:
\choiceTF
{\True Biên độ dao động của hình chiếu bằng $6$ cm.}
{\True Chu kì dao động của hình chiếu có giá trị bằng $0,4$ s.}
{\True Tần số dao động của hình chiếu bằng $2,5$ Hz.}
{Trong khoảng thời gian $1$ phút, hình chiếu thực hiện được đúng $100$ dao động toàn phần.}
\loigiai{
\begin{itemchoice}
\item (Đúng) Biên độ dao động của hình chiếu bằng $6$ cm. (Vì): Mệnh đề thứ nhất đúng vì biên độ dao động của hình chiếu bằng bán kính quỹ đạo chuyển động tròn đều: $A = R = 6$ cm
\item (Đúng) Chu kì dao động của hình chiếu có giá trị bằng $0,4$ s. (Vì): Mệnh đề thứ hai đúng vì chu kì dao động của hình chiếu bằng chu kì quay của chất điểm: $T = \dfrac{2\pi}{\omega} = \dfrac{2\pi}{5\pi} = 0,4$ s
\item (Đúng) Tần số dao động của hình chiếu bằng $2,5$ Hz. (Vì): Mệnh đề thứ ba đúng vì tần số dao động là: $f = \dfrac{1}{T} = \dfrac{1}{0,4} = 2,5$ Hz
\item (Sai) Trong khoảng thời gian $1$ phút, hình chiếu thực hiện được đúng $100$ dao động toàn phần. (Vì): Mệnh đề thứ tư sai vì số dao động thực hiện được trong $1$ phút ($\Delta t = 60$ s) là: $N = f \cdot \Delta t = 2,5 \cdot 60 = 150$ dao động
\end{itemchoice}
}
\end{ex}

% ===== Câu 27 =====
% ID: L11C1B1-10-TLN
% Mức: VD
\begin{ex}
% ID: L11C1B1-10-TLN
% Mức: VD
[SGK CTST]
	Một vật dao động điều hòa có chu kỳ $T=0,5~\mathrm{s}$. Trong thời gian $3\,\text{s}$, vật đi qua vị trí cân bằng bao nhiêu lần?
\shortans{11/01/1900}
\loigiai{
Mỗi chu kỳ vật qua VTCB 2 lần. Số chu kỳ: $3: 0,5=6\Rightarrow~\mathrm{s}$ Số lần qua VTCB= $6\cdot 2=12$.
}
\end{ex}

% ===== Câu 28 =====
% ID: L11C1B1-10-TN
% Mức: VD
\dangbt{Mối liên hệ giữa dao động điều hòa và chuyển động tròn đều. Bài toán đếm số dao động}
\begin{ex}
% ID: L11C1B1-10-TN
% Mức: VD
Trong khoảng thời gian $1$ phút, một chất điểm thực hiện được đúng $30$ dao động toàn phần quanh vị trí cân bằng. Chu kì dao động của chất điểm này là:
\choice
{\True $2$ s}
{$0,5$ s}
{$30$ s}
{$1$ s}
\loigiai{
• Phân tích các thông số thời gian và số dao động toàn phần: khoảng thời gian $\Delta t = 1 \text{ phút} = 60$ s, số dao động thực hiện được $N = 30$.
 
• Chu kì dao động $T$ là thời gian để vật thực hiện được một dao động toàn phần: $T = \dfrac{\Delta t}{N}$.
 
• Thay các giá trị vào biểu thức: $T = \dfrac{60}{30} = 2$ s.
}
\end{ex}

% ===== Câu 30 =====
% ID: L11C1B1-11-DS
% Mức: VD
\dangbt{Xác định chu kì, tần số và biên độ từ các dữ kiện quỹ đạo và thời gian}
\begin{ex}
% ID: L11C1B1-11-DS
% Mức: VD
Một vật dao động điều hòa có đồ thị li độ – thời gian dạng hình sin biến thiên tuần hoàn. Khoảng thời gian đo được giữa hai đỉnh cực đại liên tiếp trên đồ thị là $0,8$ s, và khoảng cách thẳng đứng từ đỉnh cao nhất đến đáy thấp nhất của đồ thị là $16$ cm. Đánh giá các mệnh đề sau:
\choiceTF
{Biên độ dao động của vật là $16$ cm.}
{\True Chu kì dao động của vật bằng $0,8$ s.}
{\True Tần số dao động của vật bằng $1,25$ Hz.}
{\True Tần số góc của dao động có giá trị bằng $2,5\pi$ rad/s.}
\loigiai{
\begin{itemchoice}
\item (Sai) Biên độ dao động của vật là $16$ cm. (Vì): Mệnh đề thứ nhất sai vì khoảng cách thẳng đứng từ đỉnh đến đáy đồ thị là $2A = 16$ cm, suy ra biên độ dao động của vật là $A = 8$ cm
\item (Đúng) Chu kì dao động của vật bằng $0,8$ s. (Vì): Mệnh đề thứ hai đúng vì khoảng thời gian giữa hai đỉnh cực đại (hai lần liên tiếp vật đạt biên dương) chính là một chu kì dao động $T = 0,8$ s
\item (Đúng) Tần số dao động của vật bằng $1,25$ Hz. (Vì): Mệnh đề thứ ba đúng vì tần số dao động là: $f = \dfrac{1}{T} = \dfrac{1}{0,8} = 1,25$ Hz
\item (Đúng) Tần số góc của dao động có giá trị bằng $2,5\pi$ rad/s. (Vì): Mệnh đề thứ tư đúng vì tần số góc của dao động là: $\omega = 2\pi f = 2\pi \cdot 1,25 = 2,5\pi$ rad/s
\end{itemchoice}
}
\end{ex}

% ===== Câu 31 =====
% ID: L11C1B1-11-TLN
% Mức: VD
\begin{ex}
% ID: L11C1B1-11-TLN
% Mức: VD
Một vật dao động điều hòa trên một quỹ đạo là đoạn thẳng dài $18\,\text{cm}$. Xác định biên độ dao động của vật (theo cm).
\shortans{08/01/1900}
\loigiai{
Chiều dài quỹ đạo dao động:
 $L = 2A \implies A = \dfrac{L}{2} = \dfrac{18}{2} = 9  \text{cm}.$
}
\end{ex}

% ===== Câu 32 =====
% ID: L11C1B1-11-TN
% Mức: VD
\dangbt{Mối liên hệ giữa dao động điều hòa và chuyển động tròn đều. Bài toán đếm số dao động}
\begin{ex}
% ID: L11C1B1-11-TN
% Mức: VD
Một bàn xoay hình tròn có gắn một thanh nhỏ cách tâm bàn một khoảng $15$ cm đang quay đều quanh trục với tốc độ quay đạt $2\pi$ rad/s. Khi chiếu một chùm sáng rộng, song song theo phương ngang lên một tấm màn đặt phía sau, bóng của thanh nhỏ trên màn dao động điều hòa với chu kì và tần số lần lượt là:
\choice
{\True $1$ s và $1$ Hz}
{$2$ s và $0,5$ Hz}
{$1$ s và $2$ Hz}
{$0,5$ s và $2$ Hz}
\loigiai{
• Bóng của thanh nhỏ trên màn dao động điều hòa với tần số góc bằng chính tốc độ góc quay của bàn xoay: $\omega = 2\pi$ rad/s.
 
• Chu kì dao động của bóng được xác định qua biểu thức liên hệ: $T = \dfrac{2\pi}{\omega} = \dfrac{2\pi}{2\pi} = 1$ s.
 
• Tần số dao động tương ứng của bóng trên màn là: $f = \dfrac{1}{T} = \dfrac{1}{1} = 1$ Hz.
}
\end{ex}

% ===== Câu 33 =====
% ID: L11C1B1-12-DS
% Mức: VD
\dangbt{Xác định chu kì, tần số và biên độ từ các dữ kiện quỹ đạo và thời gian}
\begin{ex}
% ID: L11C1B1-12-DS
% Mức: VD
Một vật thực hiện dao động điều hòa trên quỹ đạo là một đoạn thẳng dài $10$ cm. Biết rằng trong khoảng thời gian $10$ giây vật thực hiện được đúng $20$ dao động toàn phần. Xét các thông số động học của vật:
\choiceTF
{Biên độ dao động của vật bằng $10$ cm.}
{\True Chu kì dao động của vật bằng $0,5$ s.}
{\True Quãng đường vật đi được trong một chu kì dao động đầy đủ bằng $20$ cm.}
{\True Tốc độ trung bình của vật trong một dao động toàn phần có giá trị bằng $40$ cm/s.}
\loigiai{
\begin{itemchoice}
\item (Sai) Biên độ dao động của vật bằng $10$ cm. (Vì): Mệnh đề thứ nhất sai vì chiều dài quỹ đạo chuyển động thẳng là $L = 2A = 10$ cm, suy ra biên độ dao động của vật là $A = 5$ cm
\item (Đúng) Chu kì dao động của vật bằng $0,5$ s. (Vì): Mệnh đề thứ hai đúng vì chu kì dao động của vật là: $T = \dfrac{\Delta t}{N} = \dfrac{10}{20} = 0,5$ s
\item (Đúng) Quãng đường vật đi được trong một chu kì dao động đầy đủ bằng $20$ cm. (Vì): Mệnh đề thứ ba đúng vì quãng đường vật đi được trong một chu kì luôn bằng $S = 4A = 4 \cdot 5 = 20$ cm
\item (Đúng) Tốc độ trung bình của vật trong một dao động toàn phần có giá trị bằng $40$ cm/s. (Vì): Mệnh đề thứ tư đúng vì tốc độ trung bình được tính bằng tỉ số giữa quãng đường một chu kì và thời gian một chu kì: $v_{\text{tb}} = \dfrac{S}{T} = \dfrac{20}{0,5} = 40$ cm/s
\end{itemchoice}
}
\end{ex}

% ===== Câu 34 =====
% ID: L11C1B1-12-TN
% Mức: VD
\dangbt{Mối liên hệ giữa dao động điều hòa và chuyển động tròn đều. Bài toán đếm số dao động}
\begin{ex}
% ID: L11C1B1-12-TN
% Mức: VD
Một chất điểm dao động điều hòa. Trong $10$ dao động toàn phần liên tiếp, vật đi được tổng quãng đường dài $120$ cm. Chiều dài quỹ đạo chuyển động của chất điểm có giá trị là:
\choice
{$3$ cm}
{\True $6$ cm}
{$12$ cm}
{$9$ cm}
\loigiai{
• Trong một dao động toàn phần (một chu kì đầy đủ), quãng đường vật đi được luôn bằng $4$ lần biên độ dao động: $S_1 = 4A$.
 
• Tính biên độ từ tổng quãng đường đi được trong $10$ dao động: $S_{\text{total}} = 10 \cdot 4A = 40A$. Theo đề bài $40A = 120 \text{ cm} \Rightarrow A = 3$ cm.
 
• Chiều dài quỹ đạo chuyển động của chất điểm là khoảng cách giữa hai biên: $L = 2A = 2 \cdot 3 = 6$ cm.
}
\end{ex}

% ===== Câu 35 =====
% ID: L11C1B1-13-DS
% Mức: VD
\dangbt{Xác định chu kì, tần số và biên độ từ các dữ kiện quỹ đạo và thời gian}
\begin{ex}
% ID: L11C1B1-13-DS
% Mức: VD
Một vật dao động điều hòa thực hiện được đúng $120$ dao động toàn phần trong khoảng thời gian $1$ phút. Khảo sát các đại lượng chu kì và tần số của dao động:
\choiceTF
{\True Tần số dao động của vật bằng $2$ Hz.}
{\True Chu kì dao động của vật bằng $0,5$ s.}
{\True Tần số góc của dao động đạt giá trị $4\pi$ rad/s.}
{Khoảng thời gian ngắn nhất giữa hai lần liên tiếp vật đi qua vị trí cân bằng là $0,5$ s.}
\loigiai{
\begin{itemchoice}
\item (Đúng) Tần số dao động của vật bằng $2$ Hz. (Vì): Mệnh đề thứ nhất đúng vì khoảng thời gian $\Delta t = 1 \text{ phút} = 60$ s. Tần số dao động là số dao động thực hiện được trong một giây: $f = \dfrac{N}{\Delta t} = \dfrac{120}{60} = 2$ Hz
\item (Đúng) Chu kì dao động của vật bằng $0,5$ s. (Vì): Mệnh đề thứ hai đúng vì chu kì dao động là thời gian thực hiện một dao động: $T = \dfrac{1}{f} = \dfrac{1}{2} = 0,5$ s
\item (Đúng) Tần số góc của dao động đạt giá trị $4\pi$ rad/s. (Vì): Mệnh đề thứ ba đúng vì tần số góc của vật là: $\omega = 2\pi f = 2\pi \cdot 2 = 4\pi$ rad/s
\item (Sai) Khoảng thời gian ngắn nhất giữa hai lần liên tiếp vật đi qua vị trí cân bằng là $0,5$ s. (Vì): Mệnh đề thứ tư sai vì khoảng thời gian giữa hai lần liên tiếp vật qua vị trí cân bằng là một nửa chu kì: $\Delta t = \dfrac{T}{2} = \dfrac{0,5}{2} = 0,25$ s
\end{itemchoice}
}
\end{ex}

% ===== Câu 38 =====
% ID: L11C1B1-14-TLN
% Mức: VD
\begin{ex}
% ID: L11C1B1-14-TLN
% Mức: VD
Một vật dao động điều hòa với biên độ $8\,\text{cm}$ và chu kì dao động là $2\,\text{s}$. Tính tốc độ trung bình của vật trong một chu kì (theo cm/s).
\shortans{15/01/1900}
\loigiai{
Quãng đường trong $1$ chu kì:
 $S = 4A = 4 \cdot 8 = 32.$ v_{tb} = $\dfrac{S}{T}$ = $\dfrac{32}{2}$ = 16.
}
\end{ex}

% ===== Câu 41 =====
% ID: L11C1B1-17-TLN
% Mức: VD
\begin{ex}
% ID: L11C1B1-17-TLN
% Mức: VD
Một vật dao động điều hòa có chu kì dao động là $0,25\,\text{s}$. Tính tần số góc ($\omega$) của dao động (theo rad/s). Lấy $\pi \approx 3,14$.
\shortans{C}
\loigiai{
$\omega = \dfrac{2\pi}{T} = \dfrac{2 \cdot 3,14}{0,25} \approx 25,12.$
}
\end{ex}

% ===== Câu 43 =====
% ID: L11C1B1-18-TLN
% Mức: VD
\begin{ex}
% ID: L11C1B1-18-TLN
% Mức: VD
[SGK CTST]
	Một vật dao động điều hòa với tần số $f=2$ Hz. Tính chu kỳ dao động (kết quả làm tròn đến số nguyên phần trăm mili giây).
\shortans{500}
\loigiai{
$T=\dfrac{1}{f}=\dfrac{1}{2}=0,5\ s=500$ mili giây.
}
\end{ex}

% ===== Câu 45 =====
% ID: L11C1B1-19-TLN
% Mức: VD
\begin{ex}
% ID: L11C1B1-19-TLN
% Mức: VD
Một vật dao động điều hòa với biên độ $6\,\text{cm}$. Tính quãng đường vật đi được sau $4$ dao động toàn phần (theo cm).
\shortans{96}
\loigiai{
Trong $1$ dao động toàn phần:
 $S_1 = 4A = 4 \cdot 6 = 24.$
 Sau $4$ dao động:
 $S = 4S_1 = 4 \cdot 24 = 96.$
}
\end{ex}

% ===== Câu 47 =====
% ID: L11C1B1-20-TN
% Mức: NB
\begin{ex}
% ID: L11C1B1-20-TN
% Mức: NB
Một chất điểm dao động điều hòa trên trục Ox xung quanh gốc $O$ với biên độ $6\text{cm}$ và chu kì $2\text{s}$. Mốc để tính thời gian là khi vật đi qua vị trí $x = 3\text{cm}$ theo chiều dương. Khoảng thời gian chất điểm đi được quãng đường $249\text{ cm}$ kể từ thời điểm ban đầu là:
\choice
{\True $\dfrac{125}{6}\text{ s}$}
{$\dfrac{62}{3}\text{ s}$}
{$\dfrac{127}{6}\text{ s}$}
{$\dfrac{61}{3}\text{ s}$}
\loigiai{
Biên độ $A = 6\text{ cm}$. Chu kì $T = 2\text{ s}$.
 Thời điểm ban đầu $t=0$: vật đi qua $x=3\text{ cm}$ theo chiều dương.
 Ta có $3\text{ cm} = A/2$.
 Thời điểm ban đầu, pha ban đầu $\phi_0$: $x(0) = A\cos\phi_0 \Rightarrow 3 = 6\cos\phi_0 \Rightarrow \cos\phi_0 = 1/2$.
 Vì vật đi theo chiều dương, $v = -A\omega\sin\phi_0 {\gt } 0 \Rightarrow \sin\phi_0 {\lt } 0$.
 Vậy $\phi_0 = -\pi/3$.
 Quãng đường cần đi là $S = 249\text{ cm}$.
 Quãng đường vật đi được trong một chu kì là $4A = 4 \times 6 = 24\text{ cm}$.
 Số chu kì nguyên: $N = \left\lfloor S / (4A) \right\rfloor = \left\lfloor 249 / 24 \right\rfloor = 10$ chu kì.
 Thời gian cho $N$ chu kì là $t_N = N \times T = 10 \times 2 = 20\text{ s}$.
 Quãng đường còn lại là $\Delta S = S - N \times 4A = 249 - 10 \times 24 = 249 - 240 = 9\text{ cm}$.
 Sau $10$ chu kì, vật trở về trạng thái ban đầu: $x=3\text{ cm}$, đang đi theo chiều dương.
 Từ $x=3\text{ cm}$ (chiều dương), để đi quãng đường $\Delta S = 9\text{ cm}$:
 
 
• Vật đi từ $x=3\text{ cm}$ đến biên dương $x=A=6\text{ cm}$: quãng đường $3\text{ cm}$ ($A/2$). Thời gian $t_1 = T/6 = 2/6 = 1/3\text{ s}$.
 
• Vật đang ở $x=A=6\text{ cm}$. Quãng đường còn lại $9-3 = 6\text{ cm}$.
 
• Vật đi từ $x=A=6\text{ cm}$ về vị trí cân bằng $x=0$: quãng đường $6\text{ cm}$ ($A$). Thời gian $t_2 = T/4 = 2/4 = 1/2\text{ s}$.
 
 Tổng thời gian cho quãng đường còn lại là $\Delta t = t_1 + t_2 = 1/3 + 1/2 = 5/6\text{ s}$.
 Tổng thời gian chất điểm đi được quãng đường $249\text{ cm}$ là $t = t_N + \Delta t = 20 + 5/6 = \dfrac{125}{6}\text{ s}$.
}
\end{ex}

% ===== Câu 49 =====
% ID: L11C1B1-21-TN
% Mức: NB
\begin{ex}
% ID: L11C1B1-21-TN
% Mức: NB
Dựa vào đồ thị li độ – thời gian của một vật dao động điều hòa như hình vẽ bên, tần số góc của dao động là:
\choice
{$2\pi\ \mathrm{rad/s}$}
{\True $\pi\ \mathrm{rad/s}$}
{$10\ \mathrm{rad/s}$}
{$5\ \mathrm{rad/s}$}
\loigiai{
Chu kỳ từ đồ thị là $2\,\text{s}$, $\omega = \dfrac{2\pi}{T} = \pi\ \mathrm{rad/s}$.
}
\end{ex}

% ===== Câu 51 =====
% ID: L11C1B1-24-TN
% Mức: NB
\begin{ex}
% ID: L11C1B1-24-TN
% Mức: NB
Một vật dao động điều hoà với biên độ $4$ cm. Khi nó có li độ là $2$ cm thì vận tốc là $1$ m/s. Tần số dao động là:
\choice
{$1$ Hz}
{$3$ Hz}
{$1{,}2$ Hz}
{\True $4{,}6$ Hz}
\loigiai{
Biên độ dao động $A = 4\,\mathrm{cm}= 0,04$ m, li độ $x = 2\,\mathrm{cm}= 0,02$ m, vận tốc $v = 1$ m/s.
Ta có công thức liên hệ giữa vận tốc, li độ và tần số góc là $v = \omega \sqrt{A^2 - x^2}$.
Suy ra $\omega = \frac{v}{\sqrt{A^2 - x^2}} = \frac{1}{\sqrt{(0,04)^2 - (0,02)^2}} = \frac{1}{\sqrt{0,0016 - 0,0004}} = \frac{1}{\sqrt{0,0012}} = \frac{1}{0,02\sqrt{3}} = \frac{50}{\sqrt{3}}$ rad/s.
Tần số dao động $f = \frac{\omega}{2\pi} = \frac{50/\sqrt{3}}{2\pi} = \frac{25}{\pi\sqrt{3}}$ Hz.
Kiểm tra phương án D: $f = \frac{25}{\pi\sqrt{3}} \approx 4,6$ Hz.
Chọn D.
}
\end{ex}

% ===== Câu 53 =====
% ID: L11C1B1-25-TN
% Mức: NB
\begin{ex}
% ID: L11C1B1-25-TN
% Mức: NB
Một con ong mật đang bay tại chỗ trong không trung đập cánh với tần số khoảng $300\,\text{Hz}$. Chu kì dao động của cánh ong là
\choice
{$300\,\text{s}$}
{\True $3,33\mathrm{~ms}$}
{$3\,\text{s}$}
{$0,0033\mathrm{~s}$}
\loigiai{
Chu kì của dao động $T = \frac{1}{f} = \frac{1}{300} \approx 0,0033\mathrm{~s}$
}
\end{ex}

% ===== Câu 56 =====
% ID: L11C1B1-26-TN
% Mức: NB
\begin{ex}
% ID: L11C1B1-26-TN
% Mức: NB
Một vật nhỏ dao động điều hòa thực hiện $50$ dao động toàn phần trong $1$ s. Tần số dao động của vật là
\choice
{$50\pi ~\mathrm{Hz}$}
{$100\pi ~\mathrm{Hz}$}
{\True $50 ~\mathrm{Hz}$}
{$0,02~\mathrm{Hz}$}
\loigiai{
Tần số của dao động là số dao động thực hiện được trong $1$ giây, bằng $50$ Hz
}
\end{ex}

% ===== Câu 60 =====
% ID: L11C1B1-27-TN
% Mức: NB
\begin{ex}Một chất điểm dao động điều hòa với biên độ $A = 3$ cm và chu kỳ $T = 1$ s. Đồ thị li độ được cho như hình vẽ.
 Tính tốc độ trung bình của chất điểm trong 1 chu kỳ.
\choice
{$24\ \mathrm{cm/s}$}
{$6\ \mathrm{cm/s}$}
{\True $12\ \mathrm{cm/s}$}
{$3\ \mathrm{cm/s}$}
\loigiai{
Trong 1 chu kỳ, vật đi quãng đường $S = 4A = 12$ cm
 , $v_{tb} = \dfrac{S}{T} = \dfrac{12}{1} = 12\ \mathrm{cm/s}$
}
\end{ex}

% ===== Câu 65 =====
% ID: L11C1B1-30-TN
% Mức: TH
\begin{ex}
% ID: L11C1B1-30-TN
% Mức: TH
Một con lắc lò xo dao động với biên độ $A$. Trong một chu kì thời gian dài nhất để con lắc di chuyển từ vị trí có li độ $x_1 = -A$ đến vị trí có li độ $x_2 = A/2$ là $1\text{s}$. Chu kì dao động của con lắc là:
\choice
{$2\text{s}$}
{\True $1,5\text{s}$}
{$4\text{s}$}
{$3\text{s}$}
\loigiai{
Vật dao động điều hòa với biên độ $A$ và chu kì $T$. Vị trí ban đầu là $x_1 = -A$. Pha ban đầu là $\phi_1 = \pi$. Vị trí đích là $x_2 = A/2$. Các pha tương ứng là $\pm \pi/3$. Để tìm thời gian dài nhất từ $x_1=-A$ đến $x_2=A/2$: Ta xét góc quét trên đường tròn lượng giác từ pha $\phi_1 = \pi$ đến một trong hai pha của $x_2$.

A. Pha đích 1: $\phi_{2a} = \pi/3$. Góc quét ngắn nhất (ngược chiều kim đồng hồ) là $\Delta\alpha_1 = (2\pi - \pi) + \pi/3 = \pi + \pi/3 = 4\pi/3$. Thời gian tương ứng là $t_1 = (4\pi/3)/\omega = (4\pi/3) \times (T/2\pi) = 2T/3$.
B. Pha đích 2: $\phi_{2b} = -\pi/3$. Góc quét ngắn nhất (ngược chiều kim đồng hồ) là $\Delta\alpha_2 = \pi - (-\pi/3) = \pi + \pi/3 = 4\pi/3$. Thời gian tương ứng là $t_2 = (4\pi/3)/\omega = 2T/3$.  Cả hai trường hợp đều cho thời gian là $2T/3$. Đây là thời gian dài nhất để đi từ $x=-A$ đến $x=A/2$ (trong một chu kì). Theo đề bài, thời gian này là $1\text{ s}$. $\dfrac{2T}{3} = 1\text{ s} \Rightarrow T = \dfrac{3}{2} = 1,5\text{ s}$.
}
\end{ex}

% ===== Câu 68 =====
% ID: L11C1B1-33-TN
% Mức: NB
\begin{ex}
% ID: L11C1B1-33-TN
% Mức: NB
Một con lắc đơn gồm một hòn bi nhỏ khối lượng m, treo vào một sợi dây không giãn, khối lượng dây không đáng kể. Khi con lắc đơn này dao động điều hòa với chu kì $3\text{s}$ thì hòn bi chuyển động trên cung tròn $4\text{cm}$. Thời gian để hòn bi đi được $2\text{cm}$ kể từ vị trí cân bằng là:
\choice
{$1\text{s}$}
{\True $0,75\text{s}$}
{$2\text{s}$}
{$4\text{s}$}
\loigiai{
Chu kì $T = 3\text{ s}$. "Hòn bi chuyển động trên cung tròn $4\text{cm}$ " nghĩa là quỹ đạo dao động của hòn bi là $4\text{cm}$. Do đó, biên độ dao động $A = \dfrac{4}{2} = 2\text{ cm}$. Vật bắt đầu chuyển động từ vị trí cân bằng ($x=0$). Quãng đường cần đi là $S = 2\text{ cm}$. Vì $A=2\text{ cm}$, quãng đường $2\text{cm}$ chính là từ vị trí cân bằng đến biên ($x=0 \to x=A$ hoặc $x=0 \to x=-A$). Thời gian ngắn nhất để vật đi từ vị trí cân bằng đến biên là $t = T/4$. $t = 3/4 = 0,75\text{ s}$.
}
\end{ex}

% ===== Câu 79 =====
% ID: L11C1B1-39-TN
% Mức: NB
\begin{ex}
% ID: L11C1B1-39-TN
% Mức: NB
Một chất điểm dao động điều hòa có chu kỳ $T=1$ s. Tần số góc $\omega$ của dao động là:
\choice
{$\pi$ rad/s}
{\True $2\pi$ rad/s}
{$1$ rad/s}
{$2$ rad/s}
\loigiai{
Tần số góc $\omega$ của dao động là $\omega=\dfrac{2\pi}{T}=2\pi$ rad/s.
}
\end{ex}

\dangbt{Mối liên hệ giữa dao động điều hòa và chuyển động tròn đều. Bài toán đếm số dao động}
\begin{ex}
Một vật dao động điều hoà theo phương trình $x = 1,25\cos(2\pi t - \pi/12)$ (cm) (t đo bằng giây). Quãng đường vật đi được sau thời gian $t = 2,5$ s kể từ lúc bắt đầu dao động là
\choice
{A. $7,9$ cm.}
{B. $22,5$ cm.}
{C. $7,5$ cm.}
{\True D. $12,5$ cm.}
\loigiai{Chu kì dao động: $T = \frac{2\pi}{\omega} = \frac{2\pi}{2\pi} = 1$ s.
Thời gian dao động: $\Delta t = 2,5$ s $= 2T + 0,5T$.
Quãng đường vật đi được trong $2T$ là $S_1 = 2 \times 4A = 8A$.
Quãng đường vật đi được trong $0,5T$ là $S_2 = 2A$.
Tổng quãng đường vật đi được là $S = S_1 + S_2 = 8A + 2A = 10A = 10 \times 1,25 = 12,5$ cm.}
\end{ex}
